\section{Discussion}\label{sec:discussion}

\subsection{Data support for the narrow planet domain}

The narrow ExoEarth estimate is more model-dependent than the broad-HZ benchmark. In the DR25 population considered by \citet{Bryson2021}, only nine candidates near the HZ have $R_p\leq1.5\,\Rearth$, compared with 56 at $R_p\leq2.5\,\Rearth$. The $0.9$--$1.1\,\Rearth$ interval used here is narrower still. The result therefore depends on extrapolating a smooth fitted occurrence surface into a sparsely populated domain. This does not invalidate a model-conditioned expectation, but it limits how literally the central number should be interpreted. Theoretical and empirical studies have also warned that extrapolating from larger, more irradiated planets to small low-instellation planets can bias $\eta_\oplus$ \citep{LopezRice2018,Pascucci2019}.

The Model~2 sensitivity of +3.30\% tests one alternative temperature dependence within the Bryson framework. It does not exhaust uncertainty in the radius--instellation functional form, particularly within a domain with very few direct detections.

\subsection{Host-population uncertainty}

The JJ host count is a deterministic output of the pinned parameter realization and isochrone library. The $+0.175\%$ grid change is only a radial-discretization diagnostic and must not be read as a total uncertainty on $N_\star$. This work does not propagate posterior uncertainty in the JJ star-formation and age--metallicity histories, IMF parameters, component normalizations, or the choice among PARSEC, MIST, and BaSTI libraries supported by \texttt{jjmodel}. Consequently, the reported host count is numerically reproducible but not known astrophysically to sub-percent precision.

\subsection{Transport assumptions}\label{subsec:transport}

Equation~(\ref{eq:galactic_estimator}) transports a Kepler-calibrated occurrence model to an older Galactic population. Conditional on $\Teff$, the calculation assumes no additional dependence of planet occurrence or survival on stellar age, $[\mathrm{Fe/H}]$, $[\alpha/\mathrm{Fe}]$, thin/thick-disk membership, or Galactic radius. This is a transparent null transport assumption, not an empirical finding.

The source stellar denominator and the JJ denominator are also not guaranteed to represent identical selection or multiplicity populations. \citet{Bryson2021} began with 177,798 stars common to three source catalogs and, after cuts on stellar characterization, evolutionary state, identified binarity, noise, limb darkening, duty cycle, data span, and the DR25 \texttt{timeoutsumry} flag, retained 82,371 stars before their final temperature selections. The JJ export describes an intrinsic synthetic single-star population and does not reproduce these observational-quality cuts. The transport therefore assumes that, conditional on the modeled stellar variables, the source denominator's quality cuts are not additionally correlated with intrinsic planet occurrence.

The JJ export is constructed from a single-star IMF and single-star isochrones and does not explicitly encode a binary selection matching the Kepler cuts. Close companions can suppress planet occurrence \citep{Kraus2016}, while unresolved multiplicity can alter inferred planet radii and completeness. The direction and magnitude of a correction cannot be inferred from the present data, so multiplicity is retained as an unquantified transport systematic rather than assigned an ad hoc factor.

The empirical $g(T)$ factor is another source-population feature. Because it is embedded in the fitted occurrence model, substituting a JJ-specific $g_{\rm JJ}(T)$ after fitting would not be a controlled sensitivity experiment. A defensible replacement would require refitting the Kepler occurrence likelihood under the alternative parameterization. The Model~2 result provides a limited, internally calibrated comparison without changing the source model post hoc.

\subsection{Meaning of the Galactic annulus}

The 7--9 kpc interval is adopted solely as an absolute Galactocentric geometric mask motivated by the historically highlighted region in \citet{Lineweaver2004}. It is not a transfer of their complete GHZ model. Lineweaver et al. used $R_\odot=8.5$ kpc, whereas the pinned JJ realization uses $R_\odot=8.17$ kpc; the same absolute annulus therefore has a slightly different position relative to the Sun. The present analysis also does not inherit their time-dependent metallicity, supernova, terrestrial-planet, biological-evolution, or stellar-formation-age weighting. In particular, the one-sided criterion $t_\star\geq4.57$ Gyr is defined independently of the approximately 8--4 Gyr formation-time population emphasized in the historical GHZ result.

The concept and geometry of a sharply bounded GHZ are themselves model-dependent \citep{Prantzos2008,Gowanlock2011}. The result must therefore be described as an estimate within a 7--9 kpc annulus, not as the number of ExoEarth candidates in a physically complete Galactic habitable zone. Table~\ref{tab:domains} shows that the full JJ 4--14 kpc domain contains substantially more eligible hosts and planet intensity. These are alternative target populations, not evidence that the remainder of the disk is uninhabitable.

\subsection{What the central number does and does not mean}

The statement $\LamEE=3.38\times10^6$ means that the canonical frozen model assigns that expected number of planets to the specified radius--instellation--present-day-HZ domain around the selected host population. It is not a count of observed planets, not the number of systems with at least one such planet, and not the number of rocky, ocean-bearing, habitable, or inhabited worlds. The six-digit computational outputs retained in the archive document reproducibility; three significant figures are used in the scientific text because the model systematics are much larger.
